% TOP_PROBLEM: {"id": 20003342, "title": "Angular-factor dichotomy for triangular billiards", "category_id": 20, "category": {"id": 20, "name": "miscellaneous", "display_name": "Miscellaneous", "description": "Problems whose source classification does not fit the main mathematical categories.", "slug": "miscellaneous", "order_index": 20, "created_at": "2026-07-31T15:26:25.671Z"}, "status": "partially_solved"}
% FIRST_POSTED: 2026-09-27
% !TeX program = lualatex
% ATTEMPT_STATUS: unresolved
\documentclass[11pt]{article}
\usepackage[margin=25mm]{geometry}
\usepackage{fontspec}
\setmainfont{Latin Modern Roman}
\usepackage{amsmath,amssymb,mathtools}
\usepackage{unicode-math}
\setmathfont{Latin Modern Math}
\usepackage{xurl,hyperref}
\hypersetup{colorlinks=true,urlcolor=blue}
\setcounter{secnumdepth}{0}
\setlength{\parindent}{0pt}
\setlength{\parskip}{5pt}
\setlength{\emergencystretch}{3em}
\begin{document}
\section*{\texorpdfstring{Angular-factor dichotomy for triangular billiards}{Angular-factor dichotomy for triangular billiards}}\label{20003342}
\subsection{Definitions and mathematical statement}
Let $T\subset{\mathbb{R}}^2$ be a nondegenerate triangle with angles $\alpha_1,\alpha_2,\alpha_3$. Interpret irrational angles in the standard billiards sense:
\[
 \alpha_i/\pi\notin{\mathbb{Q}}\qquad(i=1,2,3).
\]
A unit-speed billiard moves along straight lines inside $T$ and reflects at a side with angle of incidence equal to angle of reflection. Trajectories hitting vertices form the exceptional set excluded from the almost-everywhere flow. The phase space is $T\times S^1$ with its reflection identifications and normalized Liouville measure, proportional to ${\,\mathrm{d}} x{\,\mathrm{d}}\theta$.

The question is whether every such billiard flow is ergodic: every measurable set invariant under all flow times has measure zero or one. Ergodicity of a directional component in a rational polygon is a different formulation; here the full phase space of an irrational triangle is intended.
\par\smallskip\textit{Formulation and concept references:} \href{https://aimath.org/WWN/measrigid/measrigid.pdf}{[S1]}.

\subsection{Short English statement}
Is the billiard flow in every triangle with angles irrational relative to pi ergodic on its full phase space?
\subsection{Research attempt: dense direction groups and the remaining coupled dynamics}
\textbf{Status.} Ergodicity for every triangle in the statement is not proved. We rule out all nonconstant direction-only first integrals, derive the collision-section measure and the exact reduction to its return map, and prove that regular periodic trajectories occupy a null set. These facts clarify why several plausible shortcuts do not decide full-phase ergodicity.

\textbf{Source audit.} Question 47 of [S1] is the irrational-angle ergodicity question; the preceding Question 46 concerns periodic orbits and is different. A current primary-source check also found Giovanni Forni's manuscript arXiv:2606.10102, submitted June 2026, whose abstract claims existence of a periodic orbit in every finite polygon. We do not assess that proof here. Even accepting that existence assertion would not settle Question 47: a periodic orbit has zero Liouville measure, as does the whole set of regular periodic phase points by the argument below. Historical statements about absence of a known periodic orbit must therefore not be repeated as a current ergodicity obstruction.

Write $A=\operatorname{Area}(T)$ and $P=\operatorname{Perimeter}(T)$. Interior phase points have measure
\[
 d\mu=\frac1{2\pi A}\,dx\,d\theta,
 \qquad v_\theta=(\cos\theta,\sin\theta).
\]
We work modulo the usual null sets of singular trajectories and tangent collisions. Statements about ergodicity concern measurable invariant sets modulo this measure, not whether every individual orbit is dense.

\subsubsection{1. The reflection group removes the simplest first integrals}
Let the supporting line of side $i$ have angle $\varphi_i$ modulo $\pi$. Reflection acts on velocity direction by
\[
 r_i(\theta)=2\varphi_i-\theta\pmod{2\pi}.
\]
Composing two reflections gives
\[
 r_jr_i(\theta)=\theta+2(\varphi_j-\varphi_i)\pmod{2\pi}.
                                                               \tag{1}
\]
For adjacent sides this rotation angle, up to sign and a multiple of $2\pi$, is twice their interior angle. Thus one irrational angle relative to $\pi$ already supplies an irrational circle rotation in the direction group. The assumption that all three angles are irrational is stronger than needed for this particular conclusion.

Suppose $g\in L^2(S^1)$ satisfies $g\circ r_i=g$ almost everywhere for all three sides. It is then invariant under the irrational rotation in (1). If its rotation angle is $\beta$ with $\beta/(2\pi)$ irrational, its Fourier coefficients satisfy
\[
 (e^{ik\beta}-1)\widehat g(k)=0\qquad(k\in\mathbb Z).
\]
Every coefficient with $k\ne0$ vanishes, so $g$ is constant almost everywhere. This is a measurable argument, not merely a claim about continuous invariants or topological density.

Now suppose a flow-invariant bounded function on billiard phase space has the special form $F(x,\theta)=g(\theta)$. The reflection equalities used above really do follow from flow invariance. Near a nonvertex point of a side, parameterize an incoming flow tube by side arclength, incoming angle, and time to collision. Its measure has positive density away from tangent directions. Flow invariance before and after collision implies $g(\theta)=g(r_i\theta)$ for almost every incoming angle, by Fubini on such tubes. Reflection takes incoming angles bijectively to outgoing angles; the same equality therefore holds for almost every angle on the full circle. The preceding Fourier calculation applies.

In particular, any invariant set determined solely by velocity direction has measure zero or one. There is a related elementary restriction for position-only invariants. If $F(x,\theta)=q(x)$ is invariant, free flight gives $v_\theta\cdot\nabla q=0$ distributionally for almost every $\theta$. Testing this identity against smooth compactly supported functions and integrating against $\cos\theta$ and $\sin\theta$ gives both distributional derivatives of $q$ equal to zero. Connectedness of the triangle makes $q$ constant almost everywhere. Neither conclusion rules out an invariant depending jointly on position and direction.

For comparison, a rational polygon has an explicit direction-only obstruction. If all side directions satisfy $\varphi_i-\varphi_1\in(\pi/N)\mathbb Z$, then
\[
 g(\theta)=\cos\bigl(2N(\theta-\varphi_1)\bigr)
                                                               \tag{2}
\]
is fixed by every reflection. It is nonconstant and remains unchanged during free flight, so the full billiard phase space is not ergodic. Irrationality eliminates (2), but eliminating a known obstruction is not a proof that no other obstruction exists.

\subsubsection{2. The collision section and its invariant measure}
At a side point parameterized by arclength $s$, let $\phi\in(-\pi/2,\pi/2)$ be the angle of the outgoing unit velocity from the inward normal. The collision section $\Sigma$ is the union of these side-angle rectangles, excluding endpoints and singular collision data. Write $B:\Sigma\to\Sigma$ for the next-collision map and $\tau(s,\phi)>0$ for the free-flight time.

The natural unnormalized collision measure is
\[
 d\nu_0=\cos\phi\,ds\,d\phi.                            \tag{3}
\]
To verify the factor, parameterize interior phase points after a collision by
\[
 (s,\phi,t)\longmapsto(x(s)+t v(s,\phi),\theta(s,\phi)),
 \quad0<t<\tau(s,\phi).
\]
On each straight side, the spatial Jacobian of $(s,t)$ at fixed direction is $|\det(x'(s),v)|=\cos\phi$, and $|d\theta/d\phi|=1$. Thus $dx\,d\theta=\cos\phi\,ds\,dt\,d\phi$. Specular reflection preserves angle Lebesgue measure and the magnitude of the normal velocity, so the flux measure (3) is preserved by the collision map. Equivalently, apply conservation of phase volume to a thin incoming tube and its outgoing tube.

The total mass is
\[
 \nu_0(\Sigma)=P\int_{-\pi/2}^{\pi/2}\cos\phi\,d\phi=2P.
\]
Consequently $d\nu=(2P)^{-1}d\nu_0$ is the invariant probability on the section. Almost every interior phase point has a unique immediately preceding collision. Integrating the flow-tube coordinates gives
\[
 \int_\Sigma\tau\,d\nu_0=2\pi A,
 \qquad \int_\Sigma\tau\,d\nu=\frac{\pi A}{P}.            \tag{4}
\]
The roof is integrable also directly from $0<\tau\leq\operatorname{diam}(T)$ almost everywhere. Formula (4) checks the normalization and prevents confusing uniform angle measure on the collision section with its cosine-weighted measure.

The original flow is therefore, modulo null sets, the suspension of $(\Sigma,B,\nu)$ under roof $\tau$, with probability measure
\[
 \frac{d\nu(z)\,dt}{\int\tau\,d\nu},\quad
 0\leq t<\tau(z),\qquad (z,\tau(z))\sim(Bz,0).           \tag{5}
\]
The flow is ergodic if and only if $B$ is ergodic. Indeed, a flow-invariant function is almost everywhere constant on each vertical interval, so has the form $a(z)$ there; consistency at the identifications gives $a(Bz)=a(z)$. Conversely any $B$-invariant function produces a flow-invariant function by this formula. For invariant sets the same argument applies to indicators. Since $\tau>0$ almost everywhere, a section set has zero suspension measure exactly when it has zero $\nu$ measure; complements are handled identically. This proves the equivalence in the measure class actually appearing in the problem.

The reduction is exact, but it does not say that $B$ has an expanding one-dimensional factor or a finite-state Markov description. Its continuity domains depend on which side is hit next, and their iterates develop singular boundaries. Those features must be controlled in any return-map proof.

\subsubsection{3. Regular periodic points have zero full-phase measure}
We can establish a statement stronger than the fact that one periodic orbit is null. For any fixed polygon, the set of directions of regular periodic trajectories is countable.

To prove this, specify a finite collision word $w=(i_1,\ldots,i_k)$. Unfold reflections across the corresponding side lines. Their composition is an affine Euclidean isometry
\[
 G_w(x)=L_wx+t_w,
\]
where $L_w$ is orthogonal, orientation preserving for even $k$ and reversing for odd $k$. For a trajectory returning to its initial position and velocity after this word, the unfolded straight segment has positive length $\ell$ and satisfies, in a consistent unfolding convention,
\[
 G_w(x)-x=\ell v,\qquad L_wv=v.                          \tag{6}
\]
Using the inverse convention replaces $G_w$ by its inverse and does not change the finiteness conclusion.

If $k$ is odd, $L_w$ is a reflection and has just one fixed line. It therefore permits at most two oriented unit velocities. If $k$ is even, $L_w$ is a rotation. A nonidentity rotation fixes no nonzero vector, so no periodic trajectory satisfies (6). If $L_w=I$, the first equation becomes $t_w=\ell v$. When $t_w\ne0$, this again permits at most two directions, and in fact the positive-length condition selects its orientation. If $t_w=0$, positive $\ell$ makes the equation impossible. Thus a fixed finite word permits finitely many periodic directions.

There are countably many finite words, so all regular periodic directions form a countable set $D\subset S^1$. The regular periodic phase points lie in $T\times D$, a set of $dx\,d\theta$ measure zero. This reasoning does not require rationality or irrationality of the angles.

For a fixed finite collision word ending at a vertex, unfolding similarly reduces vertex hitting to a straight segment aimed at a specified reflected vertex. For fixed $\theta$, the possible starting positions lie on a line and have area zero. Fubini and the countability of finite words give a null set of such initial states. This is the usual finite-word singularity calculation; the almost-everywhere billiard flow and its exclusion of singular trajectories are as in the statement.

Periodic orbit existence is therefore compatible with ergodicity. Nonergodicity would require an invariant set of positive measure with positive-measure complement, not a countable collection of exceptional directions or a single special orbit.

\subsubsection{4. Why dense angular motion does not solve the return-map problem}
Dense projection dynamics do not force ergodicity of a system with additional coordinates. The map
\[
 F(x,y)=(x+\alpha,y)\quad\text{on }\mathbb T^2,
\]
with irrational $\alpha$, has dense orbits in its first coordinate but preserves every set $\mathbb T\times E$. Thus its base rotation is ergodic while the full product is not. This is not a counterexample billiard: it isolates the logical error in an argument which retains only angular information.

The position-direction coupling in a billiard can be displayed explicitly. An oriented line with direction $v_\theta$ has signed transverse coordinate
\[
 b=x\cdot Jv_\theta,
\]
where $J$ rotates through $\pi/2$. This quantity is constant during straight flight. Reflection in side $i$ is an affine isometry $x\mapsto R_i x+t_i$, with $R_i$ the linear reflection. Its effect on line coordinates is
\[
 \theta'=2\varphi_i-\theta,
 \qquad b'=-b+t_i\cdot Jv_{\theta'}.                    \tag{7}
\]
Indeed, $JR_i=-R_iJ$, so
$(R_ix+t_i)\cdot JR_iv_\theta=-x\cdot Jv_\theta+t_i\cdot Jv_{\theta'}$.
The reflected line agrees with the outgoing physical line at the collision point because the affine reflection fixes that point.

An invariant can depend on the line offset $b$ as well as on $\theta$. Moreover, a reflection word is physically admissible only when the intervening segments meet the prescribed side segments in the required order. The allowed words therefore depend on the offset. The abstract density of the angular reflection group does not permit arbitrary words to be applied to every phase point while keeping all segments inside the table. Equation (7) identifies the omitted variable and the admissibility constraint concretely.

For a bounded invariant $F$, free-flight invariance locally makes $F$ a measurable function of $(\theta,b)$. The real remaining problem is to show that the side-matching relations (7), on all their admissible domains, force that joint function to be constant almost everywhere. The Fourier argument in Section 1 proves this only when the dependence on $b$ has already been removed. No argument removing it is supplied here.

\subsubsection{5. Diagnostics and conclusion of the attempt}
The finite diagnostic checks the reflection-composition formula, the affine line-offset formula (7), and the cosine normalization in (3)--(4). It also checks that a rational triangle has the explicit nonconstant invariant (2). These checks are tests of the reductions and of false shortcuts; finite orbit histograms are not used as evidence for universal ergodicity. In particular, replacing irrational angles by rational approximations changes the direction group and introduces exact invariants which the original table need not possess.

We have proved triviality of direction-only and position-only invariants, exact equivalence with the cosine-measure collision map, and nullity of all regular periodic phase points. The first unresolved step is triviality of invariants depending jointly on direction and line offset under the actual admissible reflection dynamics. Neither density of the formal reflection group nor periodic-orbit existence establishes that step. The all-irrational triangular billiard assertion remains unresolved by this attempt.

\subsection{Sources}
\begin{itemize}
\item[S1]AIM workshop problem list, Emerging Applications of Measure Rigidity.
\url{https://aimath.org/WWN/measrigid/measrigid.pdf}.
\textit{Formulation source.}
\item[S2] G. Forni, Existence of a Periodic Orbit for Billiards in Polygons, arXiv:2606.10102v1, submitted 8 June 2026. Abstract inspected; the periodic-orbit claim is not an ergodicity theorem.
\url{https://arxiv.org/abs/2606.10102v1}.
\end{itemize}
\end{document}
